\documentclass[11pt]{article}
\usepackage{amsmath,amssymb,amsthm}

\begin{document}

Maximum-norm a posteriori estimate for the backward Euler/FEM error at final time

Let \(V:=H^1_0(\Omega)\) with \(H:=L^2(\Omega)\) and identify \(H\) with its dual so that \(V\subset H\subset V^*\). Let \(c(\cdot,\cdot)\) be the elliptic bilinear form inducing a bounded linear operator \(M:V\to V^*\) by \(\langle Mv,w\rangle=c(v,w)\) for \(v,w\in V\), and set \(K:=\partial_t+M\) acting in \(V^*\). The continuous parabolic problem reads: find \(u\) with \(u(0)=u_0\), \(u|_{\partial\Omega}=0\), and such that for a.e. \(t\in(0,T)\) the variational identity holds
\[
\langle Ku(t),v\rangle=\langle F(t),v\rangle\quad\text{for all }v\in V,
\]
where the load is given by \(\langle F(t),v\rangle=(f(t),v)\) for all \(v\in V\). By the assumed regularity (in particular the Hölder continuity of \(u_0\) in space), the value \(u(T)\) is well-defined pointwise and \(\|\cdot\|_{\infty,\Omega}\) is meaningful.

Let \(0=t_0<t_1<\dots<t_N=T\) and denote by \(I_j:=(t_{j-1},t_j)\) the time slabs. Let \(u_h^j\) denote the backward Euler/FEM solution at node \(t_j\). On each slab \(I_j\) define the continuous-in-time extension \(u_h(t)\) by the prescribed interpolation rule
\[
u_h(t):=u_h^j+(t-t_j)\,\delta_t u_h^j,\qquad t\in I_j,
\]
where \(\delta_t u_h^j\) is the backward difference quotient consistent with the scheme at \(t_j\). Define the piecewise-constant load extension \(\bar f(t):=f^j\) for \(t\in I_j\) and set \(\bar F(t)\in V^*\) by \(\langle \bar F(t),v\rangle=(\bar f(t),v)\) for \(v\in V\).

Define the defect \(w:=u-u_h\). Subtracting the discrete variational identity satisfied by \(u_h\) from the continuous one for \(u\), one obtains an identity in \(V^*\) of the form
\[
\langle Kw(t),\varphi\rangle=\langle F(t)-\bar F(t),\varphi\rangle+\langle\mathcal R_{\mathrm{time}}(t),\varphi\rangle\quad\text{for all }\varphi\in V,
\]
where the first term is the load mismatch and the second term is a time-discretization residual depending only on the backward Euler scheme and on the chosen extension of \(u_h\). More precisely, because \(u_h\) is defined by \(u_h(t)=u_h^j+(t-t_j)\,\delta_t u_h^j\) on \(I_j\), the distributional time derivative \(\partial_t u_h\) on \(I_j\) is constant and equals \(\delta_t u_h^j\). Therefore the action of \(K\) on \(u_h\) is computable on each slab in the weak/distribution sense, and the above residual \(\mathcal R_{\mathrm{time}}(t)\) is determined by the discrepancy between the backward Euler discretization of \(\partial_t\) and the continuous operator \(M\) applied to the FEM spatial approximation at node values.

Next, introduce the elliptic reconstruction quantities node-wise in the standard manner for the given backward Euler/FEM method: for each node \(t_j\), define \(\psi_h^j\) and the spatial reconstruction \(R^j\) exactly as in the provided summary, so that the spatial part of the defect at time \(t_j\) is represented by a term \(R^j-u_h^j\) in the correct variational form with \(c(\cdot,\cdot)\). Concretely, the reconstruction is chosen so that for every \(\varphi\in V\),
\[
\langle M(R^j- u_h^j),\varphi\rangle=c(R^j- u_h^j,\varphi)
\]
captures precisely the spatial inconsistency produced by the FEM Galerkin projection at time \(t_j\). With this choice, the spatial inconsistencies entering \(Kw\) are exactly the reconstruction terms that can later be separated into an elliptic contribution and a time-residual contribution when paired with the Green function.

At the final time, since \(u(T)\) and \(u_h(T)\) are well-defined pointwise under the assumed regularity, we have the identity
\[
u(T)-u_h(T)=w(T).
\]
To represent \(w(T)\) by a Green function, introduce the auxiliary space-time function via the adjoint Green’s function \(\Gamma\) from the assumptions. For each \(x\in\Omega\), the function \((t\mapsto \Gamma(x,t))\) is chosen to solve the adjoint parabolic problem with homogeneous Dirichlet boundary condition \(\Gamma|_{\partial\Omega}=0\) and initial condition \(\Gamma(0)=\delta_x\) in the sense of \(V^*\) (so it is meaningful to pair \(\Gamma\) with \(V^*\)-data). The lemma’s hypotheses are designed so that the representation formula applies to any \(w\) with \(K w\) in \(V^*\); we apply it to the defect \(w\) with forcing chosen as the \(V^*\)-identity from the defect step above.

Applying Lemma \ref{lem:w-green} to \(w\) yields, for every \(x\in\Omega\),
\[
w(x,T)=(\Gamma(T),w(0))+\int_0^T\langle Kw(s),\Gamma(T-s)\rangle\,ds,
\]
where \((\Gamma(T),w(0))\) denotes the appropriate duality between \(\Gamma(T)\) and the initial defect \(w(0)=u_0-u_h(0)\). Substituting the decomposition of \(Kw(s)\) into a load mismatch part and a time-residual part, and further splitting the forcing using the node-wise elliptic reconstruction consistency encoded in \(R^j- u_h^j\), the representation becomes a sum of three types of contributions: an initial term, a time/data part involving \(\langle F-\bar F,\Gamma\rangle\) and the backward Euler residual associated to \(\delta_t u_h^j\), and a spatial reconstruction part involving \(\langle M(R^j-u_h^j),\Gamma\rangle\) (or equivalently \(c(R^j-u_h^j,\cdot)\) paired with the adjoint solution). Boundary terms vanish because \(\Gamma|_{\partial\Omega}=0\) and both \(u\) and \(u_h\) satisfy homogeneous Dirichlet boundary conditions.

Taking absolute values in the representation and using the assumed Green-function estimates---namely bounds on \(\|\Gamma(t)\|_{1,\Omega}\) and \(\|\partial_t\Gamma(t)\|_{1,\Omega}\) (equivalently on suitable \(V\)-pairs)---we obtain uniform-in-\(x\) estimates of each integral term by the corresponding dual norms of the \(V^*\)-forcing. Using the boundedness of \(c\) on \(V\times V\) and the duality \(V^*\)--\(V\), each term \(\int_0^T\langle \cdot,\Gamma(T-s)\rangle ds\) is bounded by a constant times a time-weighted integral (or sum over slabs) of the \(V^*\)-norm of the corresponding forcing component multiplied by the Green weight coming from \(\|\Gamma\|_{1,\Omega}\) (and for the time residual similarly by the weight induced through \(\|\partial_t\Gamma\|_{1,\Omega}\)). Thus, taking the supremum over \(x\in\Omega\) yields
\[
\|w(T)\|_{\infty,\Omega}\le \text{(weighted initial term)}+\text{(weighted time-residual/data mismatch term)}+\text{(weighted spatial reconstruction term)}.
\]

For the spatial reconstruction term, consider each node \(t_j\). By construction of \(R^j\) (and the associated elliptic reconstruction problem satisfied by \(y-y_h\) with right-hand side \(g\) determined by the definition of \(R^j\)), the quantity paired with the Green function at time \(t_j\) can be rewritten as an elliptic residual driven by that reconstruction. The assumed elliptic estimator property then implies that the spatial contribution at node \(j\) is bounded in maximum norm by an elliptic estimator \(\eta^j_{\mathrm{ell}}\) (with a constant independent of the mesh and time step). Therefore the sum of all spatial reconstruction contributions over nodes is bounded by a weighted sum of \(\eta^j_{\mathrm{ell}}\), where the weights are exactly the Green-function weights coming from the bound on \(\|\Gamma(\cdot)\|_{1,\Omega}\) (and, if needed, from the mapping properties used in the lemma).

For the remaining time/data contributions, the first part arises from the load mismatch \(\langle F(s)-\bar F(s),\Gamma(T-s)\rangle\). Since \(\bar F\) is piecewise constant on each \(I_j\), \(F-\bar F\) is computable from the given data and the discrete load values \(f^j\); its contribution is bounded by the corresponding dual norm \(\|F-\bar F\|_{V^*}\) weighted by the Green-function bound. The second part comes from the time-discretization residual determined solely by the backward Euler scheme, which is expressed in terms of \(\delta_t u_h^j\) and the scheme bilinear forms on each slab. By coercivity/boundedness of \(c\) (through \(M\)) and by again applying the Green estimate weights, these terms are bounded by computable time indicators depending only on discrete quantities: the values of \(\delta_t u_h^j\), the chosen time extension \(u_h(t)\) on \(I_j\), and the data mismatch \(f-\bar f\). Denote the resulting computable bound by a weighted sum of time-residual indicators \(\eta^j_{\mathrm{time}}\) (or the natural analogue arising from the \(V^*\)-norms of the backward Euler residual on each \(I_j\)).

Finally, the initial term \((\Gamma(T),w(0))\) with \(w(0)=u_0-u_h(0)\) is bounded using the Hölder continuity of \(u_0\) and the regularity/Green-function estimates given in the assumptions. This yields a nonnegative contribution that can be written in terms of the initial defect \(u_0-u_h(0)\) paired with the Green weight at time \(T\), and hence it becomes part of the final upper bound as an explicit initial-data term.

Combining all three components in the inequality for \(\|w(T)\|_{\infty,\Omega}\), using that \(u(T)-u_h(T)=w(T)\), gives the desired maximum-norm a posteriori error upper bound at final time:
\[
\|u(T)-u_h(T)\|_{\infty,\Omega}\le C\Big(\sum_{j=1}^N \omega_j\,\eta^j_{\mathrm{ell}}\;+
\sum_{j=1}^N \Omega_j\,\eta^j_{\mathrm{time}}\;+\;
\text{(initial-data term involving }u_0-u_h(0)\text{)}\Big),
\]
where \(C\) is a constant depending only on the problem data and the Green-function estimates, and where \(\omega_j,\Omega_j\) are the explicit time weights determined by the bounds in Lemma \ref{lem:w-green} (and hence by \(\|\Gamma\|_{1,\Omega}\) and \(\|\partial_t\Gamma\|_{1,\Omega}\)). All boundary contributions vanish since \(\Gamma|_{\partial\Omega}=0\) and both the primal and reconstructed functions respect the homogeneous Dirichlet setting. The elliptic contribution is isolated through \(\eta^j_{\mathrm{ell}}\), while the remaining contributions are controlled by computable time-residual/data mismatch indicators and the initial-data term, completing the proof of the final-time maximum-norm a posteriori upper bound.

\end{document}